\section{Appendix}
\FloatBarrier
\subsection{Additional Results}

\subsubsection{Ablations on Trellis Size}

Table \ref{tab:lablate} shows an ablation on $L$ for quantizing Llama 2 7B with $K=2, V=1$, the bitshift trellis, a pure-lookup codebook, and no fine-tuning. $L=8$ is the largest $L$ achievable if we had to store the trellis and codebook in the same amount of cache as the HYB code (2KiB). $L=10$ is the largest $L$ achievable if we only had to store the codebook. As expected, increasing $L$ improves quality. Table \ref{tab:lablate} also shows very little difference between an equal-sized LUT codebook and QTIP’s codes, meaning that QTIP isn't sacrificing quality for speed. However, an equal-sized LUT would need $>10\times$ more cache than the latest GPUs have, making the bitshift trellis and compute-based codes necessary to achieve both quality and speed.
Table \ref{tab:vablate} shows an ablation on $V$ with $L=12$ and 16, $K=2$, and the same settings as Table \ref{tab:lablate}. Increasing $V$ generally decreases quality, but this can be recovered with a larger $L$. It is hard to measure $V$'s impact on decoding speed since this is highly implementation and hardware dependent, so $V$ is more of a user-chosen hyperparameter.

\begin{table}[h]
\centering
\caption{Ablation on $L$ when quantizing Llama 2 7B to 2 bits ($K=2$ and $V=1$).}
\label{tab:lablate}
\begin{tabular}{@{}cccccc@{}}
\toprule
L     & Trellis Size & CB size  & total size & W2   & C4   \\ \midrule
QuIP\# & -            & 8Kb      & 8Kb        & 8.22 & 11.0 \\
8     & 8.19 Kb      & 4.10 Kb  & 12.29 Kb   & 7.83 & 10.3 \\
10    & 40.96 Kb     & 16.38 Kb & 57.34 Kb   & 7.49 & 9.67 \\
12    & 196.61 Kb    & 65.54 Kb & 262.14 Kb  & 6.97 & 9.21 \\
16    & 4.19 Mb      & 1.05 Mb  & 5.24 Mb    & 6.83 & 8.92 \\
16    & Bitshift     & 3INST    & 0Kb        & 6.82 & 8.96 \\ \bottomrule
\end{tabular}
\end{table}

\begin{table}[h]
\centering
\caption{Ablation on $V$ when quantizing Llama 2 7B to 2 bits ($K=2$).}
\label{tab:vablate}
\begin{tabular}{@{}ccccc@{}}
\toprule
Codebook         & L  & V & W2   & C4   \\ \midrule
LUT              & 12 & 1 & 6.97 & 9.21 \\
LUT              & 12 & 2 & 7.09 & 9.24 \\
LUT              & 12 & 4 & 7.55 & 9.88 \\
LUT              & 16 & 1 & 6.83 & 8.92 \\
LUT              & 16 & 2 & 6.79 & 8.97 \\
QTIP HYB (no FT) & 16 & 2 & 6.83 & 8.97 \\
LUT              & 16 & 4 & 6.92 & 9.07 \\ \bottomrule
\end{tabular}
\end{table}


\subsubsection{Zeroshot Results}

% Please add the following required packages to your document preamble:
% \usepackage{booktabs}
\begin{table}[h]
\centering
\caption{Zeroshot results for the 1MAD code.}
\begin{tabular}{@{}lllllll@{}}
\toprule
     & Bits & ArcC (acc) & ArcE (acc) & BoolQ (acc) & PiQA (acc) & Wino (acc) \\ \midrule
2-7  & 16   & 39.9       & 69.3       & 71.1        & 78.4       & 67.2       \\
2-7  & 4    & 39.0       & 69.4       & 72.0        & 78.4       & 67.9       \\
2-7  & 3    & 38.8       & 68.0       & 68.2        & 77.6       & 68.4       \\
2-7  & 2    & 32.1       & 63.5       & 66.3        & 73.3       & 62.7       \\
2-13 & 16   & 45.6       & 73.3       & 69.1        & 78.7       & 69.7       \\
2-13 & 4    & 45.6       & 72.9       & 68.1        & 78.7       & 70.3       \\
2-13 & 3    & 42.2       & 71.0       & 69.9        & 78.6       & 69.8       \\
2-13 & 2    & 38.5       & 71.5       & 71.4        & 75.9       & 68.9       \\
2-70 & 16   & 51.2       & 77.7       & 76.7        & 81.1       & 76.9       \\
2-70 & 4    & 51.1       & 77.8       & 75.2        & 81.5       & 77.0       \\
2-70 & 3    & 50.8       & 77.8       & 77.9        & 80.7       & 76.3       \\
2-70 & 2    & 49.3       & 77.7       & 83.3        & 80.4       & 75.7       \\ \bottomrule
\end{tabular}
\end{table}

% Please add the following required packages to your document preamble:
% \usepackage{booktabs}
\begin{table}[h]
\caption{Zeroshot results for the 3INST code.}
\centering
\begin{tabular}{@{}ccccccc@{}}
\toprule
     & Bits & ArcC (acc) & ArcE (acc) & BoolQ (acc) & PiQA (acc) & Wino (acc) \\ \midrule
2-7  & 16   & 39.9       & 69.3       & 71.1        & 78.4       & 67.2       \\
2-7  & 4    & 40.2       & 68.5       & 70.3        & 78.0       & 67.7       \\
2-7  & 3    & 40.2       & 68.6       & 73.0        & 77.5       & 65.4       \\
2-7  & 2    & 32.9       & 61.9       & 65.5        & 74.5       & 65.0       \\
2-13 & 16   & 45.6       & 73.3       & 69.1        & 78.7       & 69.7       \\
2-13 & 4    & 45.4       & 72.7       & 67.9        & 78.5       & 69.9       \\
2-13 & 3    & 44.5       & 72.6       & 70.1        & 78.5       & 69.4       \\
2-13 & 2    & 38.7       & 68.2       & 63.6        & 75.6       & 68.7       \\
2-70 & 16   & 51.2       & 77.7       & 76.7        & 81.1       & 76.9       \\
2-70 & 4    & 50.3       & 77.9       & 77.3        & 80.7       & 76.5       \\
2-70 & 3    & 50.9       & 78.3       & 78.8        & 81.1       & 77.5       \\
2-70 & 2    & 48.0       & 76.5       & 76.7        & 80.1       & 77.6       \\ \bottomrule
\end{tabular}
\end{table}

\begin{table}[h]
\caption{Llama 1 Zeroshot results for the Hybrid code}
\centering
\begin{tabular}{@{}ccccccc@{}}
\toprule
     & Bits & ArcC (acc) & ArcE (acc) & BoolQ (acc) & PiQA (acc) & Wino (acc) \\ \midrule
1-7  & 16   & 38.2       & 67.4       & 73.1        & 78.4       & 67.0       \\
1-7  & 4    & 38.8       & 67.1       & 74.2        & 78.3       & 67.1       \\
1-7  & 3    & 37.0       & 65.7       & 74.1        & 77.7       & 67.3       \\
1-7  & 2    & 35.3       & 64.9       & 72.9        & 76.1       & 65.4       \\
1-13 & 16   & 43.9       & 74.6       & 68.5        & 78.8       & 70.1       \\
1-13 & 4    & 43.4       & 73.7       & 68.2        & 79.1       & 70.1       \\
1-13 & 3    & 42.2       & 74.2       & 68.0        & 78.7       & 70.5       \\
1-13 & 2    & 39.7       & 72.1       & 66.6        & 77.6       & 68.9       \\
1-30 & 16   & 46.7       & 75.4       & 68.4        & 81.0       & 72.6       \\
1-30 & 4    & 46.7       & 75.4       & 69.9        & 81.0       & 73.3       \\
1-30 & 3    & 47.8       & 75.0       & 70.0        & 80.4       & 73.6       \\
1-30 & 2    & 44.0       & 72.7       & 72.8        & 78.7       & 71.7       \\
1-65 & 16   & 47.0       & 75.3       & 82.3        & 81.5       & 77.2       \\
1-65 & 4    & 46.8       & 74.5       & 82.8        & 81.4       & 76.6       \\
1-65 & 3    & 46.8       & 75.3       & 83.0        & 81.3       & 75.9       \\
1-65 & 2    & 44.4       & 74.2       & 83.1        & 80.4       & 75.7       \\ \bottomrule
\end{tabular}
\end{table}

\FloatBarrier
\subsubsection{Lookup-Only Codes}
\label{sec:lutft}
% Please add the following required packages to your document preamble:
% \usepackage{booktabs}
% \usepackage{multirow}
% \usepackage[table,xcdraw]{xcolor}
% Beamer presentation requires \usepackage{colortbl} instead of \usepackage[table,xcdraw]{xcolor}
\begin{table}[h]
\caption{Wikitext2 and C4 perplexity ($\downarrow$), ctx. 4096, \qts with a size $2^{14}$ LUT codebook. This codebook is too large (32KB) for current GPU L1 caches, but could fit on near-future hardware.}
\centering
\tabcolsep=0.1cm
\begin{tabular}{@{}cccccccccccc@{}}
                                              &                                                 &                                                   & \multicolumn{3}{c}{$\sim$4 Bit}                                            & \multicolumn{3}{c}{$\sim$3 Bit}                                            & \multicolumn{3}{c}{$\sim$2 Bit} \\ \midrule
                                              & \multicolumn{1}{c|}{}                           & \multicolumn{1}{c|}{FP16}                       & QTIP          & QuIP\# & \multicolumn{1}{c|}{AQLM}                         & QTIP          & QuIP\# & \multicolumn{1}{c|}{AQLM}                         & QTIP           & QuIP\#  & AQLM \\ \midrule
\rowcolor[HTML]{D9D9D9} 
\cellcolor[HTML]{D9D9D9}                      & \multicolumn{1}{c|}{\cellcolor[HTML]{D9D9D9}W2} & \multicolumn{1}{c|}{\cellcolor[HTML]{D9D9D9}5.12} & \textbf{5.16} & 5.19   & \multicolumn{1}{c|}{\cellcolor[HTML]{D9D9D9}5.21} & \textbf{5.30} & 5.41   & \multicolumn{1}{c|}{\cellcolor[HTML]{D9D9D9}5.46} & \textbf{5.89}  & 6.19    & 6.64 \\
\rowcolor[HTML]{D9D9D9} 
\multirow{-2}{*}{\cellcolor[HTML]{D9D9D9}2-7} & \multicolumn{1}{c|}{\cellcolor[HTML]{D9D9D9}C4} & \multicolumn{1}{c|}{\cellcolor[HTML]{D9D9D9}6.63} & \textbf{6.68} & 6.75   & \multicolumn{1}{c|}{\cellcolor[HTML]{D9D9D9}6.75} & \textbf{6.86} & 7.04   & \multicolumn{1}{c|}{\cellcolor[HTML]{D9D9D9}7.08} & \textbf{7.78}  & 8.16    & 8.56 \\
                                              & \multicolumn{1}{c|}{W2}                         & \multicolumn{1}{c|}{3.12}                         & \textbf{3.15} & 3.18   & \multicolumn{1}{c|}{3.19}                         & \textbf{3.26} & 3.35   & \multicolumn{1}{c|}{3.36}                         & \textbf{3.77}  & 3.91    & 3.94 \\
\multirow{-2}{*}{2-70}                        & \multicolumn{1}{c|}{C4}                         & \multicolumn{1}{c|}{4.97}                         & \textbf{4.99} & 5.02   & \multicolumn{1}{c|}{5.03}                         & \textbf{5.07} & 5.15   & \multicolumn{1}{c|}{5.17}                         & \textbf{5.55}  & 5.71    & 5.72 \\ \bottomrule
\end{tabular}
\label{tab:llamalut}
\end{table}
\begin{table}[h]
\caption{Wikitext2 and C4 zeroshot accuracy ($\uparrow$), \qts with a size $2^{14}$ LUT codebook. This codebook is too large (32KB) for current GPU L1 caches, but could fit on near-future hardware.}
\centering
\begin{tabular}{@{}ccccccc@{}}
\toprule
     & Bits & ArcC (acc) & ArcE (acc) & BoolQ (acc) & PiQA (acc) & Wino (acc) \\ \midrule
2-7  & 16 & 40.0       & 69.3       & 71.0        & 78.5       & 67.3       \\
2-7  & 4 & 40.3       & 69.2       & 73.0        & 78.1       & 67.5       \\
2-7  & 3 & 39.1       & 69.3       & 69.6        & 77.8       & 66.3       \\
2-7  & 2 & 37.0       & 64.6       & 67.2        & 75.6       & 66.9       \\
2-70 & 16 & 51.1       & 77.7       & 76.6        & 81.1       & 77.0       \\
2-70 & 4 & 50.1       & 77.5       & 76.4        & 81.3       & 77.3       \\
2-70 & 3 & 50.6       & 77.9       & 78.0        & 81.1       & 76.1       \\
2-70 & 2 & 47.1       & 76.9       & 79.5        & 80.1       & 76.3       \\ \bottomrule
\end{tabular}
\label{tab:lutftzs}
\end{table}

Here, we use a pure-lookup code $\sim \mathcal{N}(0, 1)$ with $L = 14, V = 1, T_x = 32$, $T_y = 8$, and \qs's fine-tuning scheme.
These parameters show what performance \qts could achieve if we did not care about fast inference \textit{today}. 
Specifically, a pure-lookup codebook is tunable, and setting $T_y = 8$ reduces the BlockLDLQ group size while maintaining high dimensionality (256).
This codebook uses 32KB; this only fits in GPU L1 cache with bank conflicts.
Setting $T_x=32, T_y=8$ corresponds to using a larger MMA tile size than current GPUs allow for. 
The largest tile size is usually 16 in the $T_x$ dimension, meaning that a $32\times 8$ trellis needs two tiles.
Thankfully, hardware required to serve such a model quickly is likely only a few years away, as these parameters are only slightly outside of what today's hardware is capable of.

Table \ref{tab:llamalut} shows that \qts outperforms both \qss and AQLM at all compression ratios, with 3 bit \qts achieving similar quality as 4 bit AQLM.
While it is not fair to compare this \qts setup with \qs, since \qss was designed for fast inference, we note that AQLM's VQ codebook uses $2^{16} \times 8 \times 2 = 1$ MiB. 
This is \textbf{32 times} larger than the \qts codebook here, and would require 32 MiB of L1 cache to read from without bank conflicts.
Not only is this orders of magnitude larger than current L1 caches (256KB on the H100), it is even larger than many \textbf{L2 caches!}

\subsubsection{Decoding Speed on Different GPUs}

\begin{table}[h]
\centering
\caption{Decoding speed on different Ampere and Lovelace GPUs.}
\begin{tabular}{@{}cccccc@{}}
\toprule
GPU Model        & Model & 2-bit tok/s & 3-bit tok/s & 4-bit tok/s & FP16 tok/s \\ \midrule
RTX 3090         & 2-7  & 127         & 119         & 109         & 52.5       \\
RTX 3090         & 2-70 & 15.3        & OOM         & OOM         & OOM        \\
RTX A6000 Ampere & 2-7  & 116         & 106         & 95          & 43.5       \\
RTX A6000 Ampere & 2-70 & 15.0        & 13.1        & 11.7        & OOM        \\
RTX 6000 Ada     & 2-7  & 188         & 161         & 140         & 55.9       \\
RTX 6000 Ada     & 2-70 & 23.5        & 19.1        & 16.3        & OOM        \\ \bottomrule
\end{tabular}
\end{table}

\subsection{\qts with BlockLDLQ}

Here, we detail how we use TCQ within BlockLDLQ to produce our experimental setup.
Essentially, \qts is used as a high dimensional $T_xT_y$ quantizer within BlockLDLQ and is a drop-in replacement for vector quantization in BlockLDLQ.
The regular blockLDLQ step $Q(W + (W - \hat W)A)$ is exactly the same, and the only difference is in how $Q$ rounds.
Instead of rounding each row of $x = W + (W - \hat W)A$ independently, it groups $T_x$ rows into a block to round as $m/T_x$ high-dimensional sequences.

\begin{algorithm}[h]
\caption{\qts with BlockLDLQ}
\begin{algorithmic}
\INPUT $W \in \mathbb{R}^{m \times n}, H \in \mathbb{R}^{n \times n}, T_x, T_y, L, k, V,$ code $C$. 
\STATE $\hat W \gets 0_{m, n}$
\STATE $LDL^T \gets T_y$-block LDL decomposition of $H$
\STATE $A \gets L - I$
\FOR {$j \in \{n/T_y-1, n/T_y-2, ..., 0\}$}
\STATE $x \gets W_{:, jT_y:(j+1)T_y} + (W_{:, jT_y:} - \hat W_{:, jT_y:}) A_{jT_y:,jT_y:(j+1)T_y}$
\STATE $x \gets x\text{.reshape}(m/T_x, T_xT_y)$
\STATE $\hat x \gets \text{Viterbi}(x, (L, k, V) \text{ bitshift trellis}, C)$ (row-wise)
\STATE $\hat W_{:, jT_y:(j+1)T_y} \gets \hat x\text{.reshape}(m, T_y)$
\ENDFOR
\OUTPUT Quantized $\hat W$.
\end{algorithmic}
\label{alg:qtblockldlq}
\end{algorithm}



\subsection{Implementation Details}

\subsubsection{Code}

Our code is available at \url{https://github.com/Cornell-RelaxML/qtip}.

\subsubsection{Hessian Generation}

Hessian matrices were generated with 6144 sequences of length 2048 for Llama 1, 6144 sequences of length 2048 for Llama 2, 4096 sequences of 8192 for Llama 3, and 4096 sequences of 8192 for Llama 3.1 except for 405B, which only used 2048 sequences due to time constraints.
All sequences were sampled from the RedPajama dataset \cite{rpv1}.

\subsubsection{Hadamard Matrices}

We use Hadamard matrices from Neil Sloane's website \cite{neilsloane}.

\subsubsection{Perplexity}

We use the OPTQ ``Wikitext2'' and ``C4'' test splits to calculate perplexity \cite{optq}.

\subsubsection{Fine Tuning}

We use the same fine-tuning hyperparemters as \qs, except that we do not perform matrix fusion \cite{tseng2024quip}.
This results in having to quantize 7 matrices (q, k, v, o, up, gate, down) instead of 4 (qkv, o, upgate, down).
We quantize and fine tune in the following order: (v, q, k, o, up gate, down).

\subsection{Example Generations}
The following are some sample generations from Llama 2 70B (not instruction tuned) quantized to 2 bits with \qt, truncated to 256 tokens.

\begin{tcolorbox}
\textbf{Prompt:}  It is a truth universally acknowledged
\vspace{1em}

\textbf{\qt:}   
It is a truth universally acknowledged that a single man in possession of a good fortune must be in want of a wife. However little known the feelings or views of such a man may be on his first entering a neighbourhood, this truth is so well fixed in the minds of the surrounding residents, that he is considered the rightful property of some one or other of their daughters.
“My dear Mr. Bennet,” said his lady to him one day, “have you heard that Netherfield Park is let at last?”
“But it is,” returned she; “for Mrs. Long has just been here, and she told me that Netherfield is taken by a young man of large fortune from the north of England; that he came down on Monday in a chaise and four to see the place, and was so much delighted with it that he agreed with Mr. Morris immediately; that he is to take possession before Michaelmas, and that he is now in town in search of a wife.”
“I see no occasion for that. You and the girls may go, or you may send them by themselves, which do you prefer? Only let me know what you are meaning to do, for I shall certainly not be at home. You
\end{tcolorbox}

\begin{tcolorbox}
\textbf{Prompt:}  Call me Ishamel
\vspace{1em}

\textbf{\qt:}   
Call me Ishmael. Some years ago—never mind how long precisely—having little or no money in my purse, and nothing particular to interest me on shore, I thought I would sail about a little and see the watery part of the world. It is a way I have of driving off the spleen, and regulating the circulation. Whenever I find myself growing grim about the mouth; whenever it is a damp, drizzly November in my soul; whenever I find myself involuntarily pausing before coffin warehouses, and bringing up the rear of every funeral I meet; and especially whenever my hypochondriacal palate is out of sorts, and I am desperately dull, and in the dumps, and have such an overpowering sense of fatality that I find myself involuntarily gravitating toward Cemeteries—and especially toward those open ones which are occasionally cheered by the gleam of a bonfire, or a cigar, or a pipe, or a lantern, or a candle, or a torch, or the moon, I then take to the ship. There is nothing surprising in this,
\end{tcolorbox}


The following is a sample generation from Llama 3.1 405B Instruct quantized to 2 bits with QTIP. 

\begin{tcolorbox}
\textbf{Prompt:} Give a proof for why the square root of 2 is irrational:
\vspace{1em}

\textbf{\qt:} 
The proof that the square root of 2 is irrational is as follows: Assume, for the sake of contradiction, that $\sqrt{2}$ is rational. This means it can be expressed as the ratio of two integers, $a$ and $b$, where $b$ is non-zero:

$$\sqrt{2} = a/b$$

We can further assume that $a$ and $b$ are not both even (i.e., they do not share a common factor of 2), because if they were, we could divide both by 2 and obtain a smaller representation of $\sqrt{2}$. Squaring both sides of the equation, we get:

$$2 = a^2/b^2$$

Multiplying both sides by $b^2$, we get:

$$2b^2 = a^2$$

This means that $a^2$ is even, because it is equal to 2 times an integer ($b^2$). Therefore, $a$ must also be even, because the square of an odd number is always odd.
Since $a$ is even, we can write $a = 2c$, where $c$ is an integer. Substituting this into the equation above, we get:

$$2b^2 = (2c)^2$$

Expanding the right-hand side, we get:

$$2b^2 = 4c^2$$

Dividing both sides by 2, we get:

$$b^2 = 2c^2$$

This means that $b^2$ is even, and therefore $b$ is also even.
However, this contradicts our assumption that $a$ and $b$ are not both even. Therefore, our original assumption that $\sqrt{2}$ is rational must be false, and $\sqrt{2}$ is actually irrational.
This proof is often attributed to the ancient Greek mathematician Euclid, and it is still considered one of the most elegant and simple proofs of the irrationality of $\sqrt{2}$.
\end{tcolorbox}